\documentclass[a4paper,12pt,oneside]{article}

\usepackage[T2A]{fontenc}
\usepackage[utf8]{inputenc}
\usepackage[russian]{babel}

\usepackage{tempora}
\usepackage{microtype}

\usepackage{amsmath,amssymb,amsthm,mathtools}

\usepackage{geometry}
\geometry{top=25mm,left=25mm,right=20mm,bottom=25mm}

\usepackage{enumitem}
\usepackage{graphicx}
\setlist[enumerate]{leftmargin=*, itemsep=0.55\baselineskip, topsep=0.35\baselineskip}
\setlist[itemize]{leftmargin=*, itemsep=0.25\baselineskip, topsep=0.25\baselineskip}

\setlength{\parindent}{1.25em}
\setlength{\parskip}{0pt}

\newcommand{\DefItem}[2]{\item \textbf{#1}\\#2}

\begin{document}

\begin{center}
    {\Large\bfseries Домашняя работа №1. Теория множеств.}\\[2mm]
\end{center}
\vspace{2mm}
\hrule
\vspace{4mm}

\section*{Задание 0}
Определите истинность заданных утверждений. Считайте, что $a \ne b$.
\begin{enumerate}[label=\arabic*.]
\item $a \in 2^{\{a\}}$.
\item $2^{\{a,\varnothing\}} \subset 2^{\{a,b,\varnothing\}}$.
\item $\{a,b\} \subseteq 2^{\{a,b\}}$.
\end{enumerate}

\section*{Задание 1}
\[
A=\{0,3,5,6\},\quad
B=\{0,1,2,3,5,7\},\quad
C=\{1,4,5,6\},\quad
D=\{0,6,7,8\},\quad
U=\{0,1,2,3,\ldots,9\}.
\]
Нарисуйте диаграмму Венна и найдите элементы множества
\[
P=(\overline{A}\cap \overline{B})
\cup (B\cap C\cap \overline{D})
\cup (B\cap \overline{C}\cap \overline{D})
\cup (A\cap \overline{B}\cap D).
\]

\section*{Задание 2}
$U$ – универсум. Упростите выражения, если $A\subset B$, $B=C$.
\begin{enumerate}[label=\arabic*.]
\item $\overline{A\cup \overline{\overline{B}\cap C}}$.
\item $\overline{A\cap \overline{B}\cap C}\ \cup\ U$.
\end{enumerate}

\section*{Задание 3}
Упростить выражение, используя законы склеивания, поглощения, де Моргана и т.\,п.
\begin{enumerate}[label=\arabic*.]
\item $(A\cup \overline{B})\cap (A\cup \overline{B}\cup C)\cap (A\cup \overline{B}\cup D)$.
\item $\overline{A}\cap \overline{B}\ \cup\ A\cap B\ \cup\ \overline{A}\cap B$.
\item $\overline{A\cup A\cap B\cup A\cap C}$.
\end{enumerate}

\section*{Задание 4}
$U$ – универсум. Упростить формулу при условии, что множества $A,B,C,D$ связаны отношением
\[
A\subset B\subset C\subset D\subset U.
\]
\[
P=\overline{A}\cap B\cap \overline{C}\ \cup\ \overline{A}\cap C\cap D\ \cup\ A\cap B\cap \overline{C}\ \cup\ A\cap B\cap C.
\]


\section*{Задание 5}
Воспользовавшись свойствами (законами) операций над множествами, выяснить, верны ли тождества для любых $A,B,C$.
\begin{enumerate}[label=\arabic*.]
\item $A \cap (B \Delta C) = (A \cap B)\Delta (A \cap C)$.
\item $A \setminus (A \setminus B) = A \cap B$.
\end{enumerate}

\section*{Задание 6}
Задайте множество аналитически по закрашенной области на диаграмме Эйлера–Венна. Упростите полученные выражения.

\[
\includegraphics[width=0.5\linewidth]{HW1.diagram.png}
\]


\section*{Задание 7}
\begin{enumerate}[label=\arabic*.]
\item Найдите множества $A$, $B$ и $C$, если
\[
A\times B=\{(1,4),(2,3),(1,3),(2,4)\},
\]
\[
B\times C=\{(3,5),(4,5),(4,6),(4,7),(3,7),(3,6)\}.
\]
\item Пусть множество $A$ состоит из $6$ элементов, множество $B$ – из $4$ элементов. При этом $|A\cap B|=3$. Найдите
\[
|A\times B\times (A\cup B)|.
\]
\item Найдите мощность множества
\[
(A\times (B\cap A))\cup (A\times B),
\]
если
\[
A=\{m,n,k,l\},\qquad B=\{a,b,c,i,j,k\}.
\]
\end{enumerate}

\section*{Задание 8}
\begin{enumerate}[label=\arabic*.]
\item Пусть $A=\{3,|B|\}$, $B=\{1,|A|,|B|\}$. Найдите, чему равны множества $A$ и $B$.
\item Даны множества
\[
A=\{a,b,c,d\},\quad
B=\{a,b,4\},\quad
C=\{4,2,c\},\quad
D=\{a,b,3\},\quad
E=\{1,b\}.
\]
Найдите $a,b,c,d$, зная, что $B\subseteq A$, $C\subseteq A$, $D\subseteq A$, $E\subseteq B$.
\end{enumerate}

\section*{Задание 9}
Предположим, что $|A|=m$ и $|B|=n$. Найдите мощность множества
\begin{enumerate}[label=\arabic*.]
\item $\bigl|\{\,X\subseteq P(A)\mid |X|\le 1\,\}\bigr|$.
\end{enumerate}

\end{document}
